\documentclass[10pt,a4paper]{amsart}
\usepackage{amsmath,amssymb,amsthm,graphicx}
\usepackage{verbatim}
\newenvironment{reference}{\comment}{\endcomment}
%\newenvironment{reference}{}{}
\newenvironment{slogan}{\comment}{\endcomment}
\newenvironment{history}{\comment}{\endcomment}
% For commutative diagrams we use Xy-pic
\usepackage[all]{xy}

% We use 2cell for 2-commutative diagrams.
\xyoption{2cell}
% Theorem environments.
%
\theoremstyle{plain}
\newtheorem{theorem}[subsection]{Theorem}
\newtheorem{proposition}[subsection]{Proposition}
\newtheorem{lemma}[subsection]{Lemma}

\theoremstyle{definition}
\newtheorem{definition}[subsection]{Definition}
\newtheorem{example}[subsection]{Example}
\newtheorem{exercise}[subsection]{Exercise}
\newtheorem{situation}[subsection]{Situation}

\theoremstyle{remark}
\newtheorem{remark}[subsection]{Remark}
\newtheorem{remarks}[subsection]{Remarks}

\numberwithin{equation}{subsection}

% Macros
%
\def\lim{\mathop{\mathrm{lim}}\nolimits}
\def\colim{\mathop{\mathrm{colim}}\nolimits}
\def\Spec{\mathop{\mathrm{Spec}}}
\def\Hom{\mathop{\mathrm{Hom}}\nolimits}
\def\Ext{\mathop{\mathrm{Ext}}\nolimits}
\def\SheafHom{\mathop{\mathcal{H}\!\mathit{om}}\nolimits}
\def\SheafExt{\mathop{\mathcal{E}\!\mathit{xt}}\nolimits}
\def\Sch{\mathit{Sch}}
\def\Mor{\mathop{\mathrm{Mor}}\nolimits}
\def\Ob{\mathop{\mathrm{Ob}}\nolimits}
\def\Sh{\mathop{\mathit{Sh}}\nolimits}
\def\NL{\mathop{N\!L}\nolimits}
\def\CH{\mathop{\mathrm{CH}}\nolimits}
\def\proetale{{pro\text{-}\acute{e}tale}}
\def\etale{{\acute{e}tale}}
\def\QCoh{\mathit{QCoh}}
\def\Ker{\mathop{\mathrm{Ker}}}
\def\Im{\mathop{\mathrm{Im}}}
\def\Coker{\mathop{\mathrm{Coker}}}
\def\Coim{\mathop{\mathrm{Coim}}}

% Boxtimes
%
\DeclareMathSymbol{\boxtimes}{\mathbin}{AMSa}{"02}

%
% Macros for moduli stacks/spaces
%
\def\QCohstack{\mathcal{QC}\!\mathit{oh}}
\def\Cohstack{\mathcal{C}\!\mathit{oh}}
\def\Spacesstack{\mathcal{S}\!\mathit{paces}}
\def\Quotfunctor{\mathrm{Quot}}
\def\Hilbfunctor{\mathrm{Hilb}}
\def\Curvesstack{\mathcal{C}\!\mathit{urves}}
\def\Polarizedstack{\mathcal{P}\!\mathit{olarized}}
\def\Complexesstack{\mathcal{C}\!\mathit{omplexes}}
% \Pic is the operator that assigns to X its picard group, usage \Pic(X)
% \Picardstack_{X/B} denotes the Picard stack of X over B
% \Picardfunctor_{X/B} denotes the Picard functor of X over B
\def\Pic{\mathop{\mathrm{Pic}}\nolimits}
\def\Picardstack{\mathcal{P}\!\mathit{ic}}
\def\Picardfunctor{\mathrm{Pic}}
\def\Deformationcategory{\mathcal{D}\!\mathit{ef}}
\usepackage{cleveref}
\crefname{originalreference}{original reference}{original references}
\Crefname{originalreference}{Original reference}{Original references}
\makeatletter
\newlabel{spaces-morphisms-lemma-quasi-compact-local}{{spaces-morphisms-lemma-quasi-compact-local}{source}{}{sourcepointer-1}{}}
\newlabel{spaces-morphisms-lemma-quasi-compact-local@cref}{{[originalreference][0][]spaces-morphisms-lemma-quasi-compact-local}{[1][0][]source}{}{}{}}
\makeatother
\crefname{originalreference}{original reference}{original references}
\Crefname{originalreference}{Original reference}{Original references}
\makeatletter
\newlabel{stacks-properties-lemma-quasi-compact-stack}{{stacks-properties-lemma-quasi-compact-stack}{source}{}{sourcepointer-2}{}}
\newlabel{stacks-properties-lemma-quasi-compact-stack@cref}{{[originalreference][0][]stacks-properties-lemma-quasi-compact-stack}{[1][0][]source}{}{}{}}
\makeatother
\crefname{originalreference}{original reference}{original references}
\Crefname{originalreference}{Original reference}{Original references}
\makeatletter
\newlabel{stacks-properties-section-properties-morphisms}{{stacks-properties-section-properties-morphisms}{source}{}{sourcepointer-3}{}}
\newlabel{stacks-properties-section-properties-morphisms@cref}{{[originalreference][0][]stacks-properties-section-properties-morphisms}{[1][0][]source}{}{}{}}
\makeatother
\title{For a representable morphism, quasi-compactness is tested after arbitrary base change}
\author{The Stacks Project Authors}
\date{stacks-project:stacks-morphisms.tex:local-snapshot-20261012; GNU FDL 1.2 or later}
\begin{document}
\maketitle
\noindent\textit{Editorial note.} All delivered mathematical text is the author's own, copied from the pinned Stacks Project chapter source sources/stacks-algebraic-stacks-morphisms/pinned-original-source.tex; no statement or proof step was rewritten, reordered or removed.\par\medskip
\noindent\textit{Editorial note.} Preamble provenance: the Stacks Project keeps one single shared preamble for every chapter; this card takes it directly from the staged sources/stacks-algebraic-stacks-morphisms/preamble.tex as three verbatim spans (lines 8-12, 14-18 and 154-217), so no derived file was written and nothing was added or redefined. The wrapper class and the base packages are supplied by the assembler.\par\medskip
\noindent\textit{Editorial note.} Bibliography provenance: neither the counted statement nor its proof contains a citation (this lane rescans the whole span before writing the card), so the card carries no bibliography and no entry was reconstructed.\par\medskip
\noindent\textit{Editorial note.} Selection provenance: this card was opened under the file-based control contract (CONTROL.md), with the mechanical pre-card gate run on its statement and proof interval before the card was written.\par\medskip
\noindent\textit{Allowed standard prerequisites (editorial).}\begin{itemize}
\item The description of a quasi-compact algebraic stack through a surjective smooth morphism from a quasi-compact scheme. Invoked twice inside the counted proof (source lines 1935-1937 and 1947-1950) as a cross-chapter pointer. Supplies the presentation of a quasi-compact test object and descends quasi-compactness along a smooth surjection.
\item The local criterion for quasi-compactness of a morphism of algebraic spaces. Invoked at the end of the counted proof (source lines 1958-1960) as a cross-chapter pointer. Converts quasi-compactness of the preimages of affine opens into quasi-compactness of the morphism.
\item The definition of quasi-compactness for a representable morphism. Used inside the counted proof (source lines 1938-1940). Gives that the induced morphism of algebraic spaces is quasi-compact.
\end{itemize}\medskip
\ifdefined\hypertarget\hypertarget{sourcepointer-1}{}\fi
\noindent\textit{Source reference.} Original reference spaces-morphisms-lemma-quasi-compact-local: sources/stacks-algebraic-stacks-morphisms/pinned-original-source.tex, source lines 1960--1960. Navigation pointer only: the counted text refers to the Stacks item spaces-morphisms-lemma-quasi-compact-local, which belongs to another chapter of the project and is not part of this card. The pointer carries the target's own label as its number because this chapter's staged snapshot does not contain the other chapter's tag table; its content is not reproduced here.\par\medskip
\ifdefined\hypertarget\hypertarget{sourcepointer-2}{}\fi
\noindent\textit{Source reference.} Original reference stacks-properties-lemma-quasi-compact-stack: sources/stacks-algebraic-stacks-morphisms/pinned-original-source.tex, source lines 1935--1935. Navigation pointer only: the counted text refers to the Stacks item stacks-properties-lemma-quasi-compact-stack, which belongs to another chapter of the project and is not part of this card. The pointer carries the target's own label as its number because this chapter's staged snapshot does not contain the other chapter's tag table; its content is not reproduced here.\par\medskip
\ifdefined\hypertarget\hypertarget{sourcepointer-3}{}\fi
\noindent\textit{Source reference.} Original reference stacks-properties-section-properties-morphisms: sources/stacks-algebraic-stacks-morphisms/pinned-original-source.tex, source lines 248--248. Navigation pointer only: the counted text refers to the Stacks item stacks-properties-section-properties-morphisms, which belongs to another chapter of the project and is not part of this card. The pointer carries the target's own label as its number because this chapter's staged snapshot does not contain the other chapter's tag table; its content is not reproduced here.\par\medskip
\begin{lemma}
\label{lemma-characterize-representable-quasi-compact}
Let $f : \mathcal{X} \to \mathcal{Y}$ be a morphism of algebraic stacks
which is representable by algebraic spaces. The following are equivalent:
\begin{enumerate}
\item $f$ is quasi-compact (as in Properties of Stacks,
Section \ref{stacks-properties-section-properties-morphisms}), and
\item for every quasi-compact algebraic stack $\mathcal{Z}$
and any morphism $\mathcal{Z} \to \mathcal{Y}$ the algebraic stack
$\mathcal{Z} \times_\mathcal{Y} \mathcal{X}$ is quasi-compact.
\end{enumerate}
\end{lemma}

\begin{proof}
Assume (1), and let $\mathcal{Z} \to \mathcal{Y}$
be a morphism of algebraic stacks with $\mathcal{Z}$ quasi-compact. By
Properties of Stacks,
Lemma \ref{stacks-properties-lemma-quasi-compact-stack}
there exists a quasi-compact scheme $U$ and a surjective smooth
morphism $U \to \mathcal{Z}$. Since $f$ is representable by algebraic
spaces and quasi-compact we see by definition that
$U \times_\mathcal{Y} \mathcal{X}$ is an algebraic space, and that
$U \times_\mathcal{Y} \mathcal{X} \to U$ is quasi-compact.
Hence $U \times_\mathcal{Y} \mathcal{X}$ is a quasi-compact algebraic space.
The morphism
$U \times_\mathcal{Y} \mathcal{X} \to
\mathcal{Z} \times_\mathcal{Y} \mathcal{X}$
is smooth and surjective (as the base change of the smooth
and surjective morphism $U \to \mathcal{Z}$).
Hence $\mathcal{Z} \times_\mathcal{Y} \mathcal{X}$
is quasi-compact by another application of
Properties of Stacks,
Lemma \ref{stacks-properties-lemma-quasi-compact-stack}

\medskip\noindent
Assume (2). Let $Z \to \mathcal{Y}$ be a morphism, where $Z$ is a scheme.
We have to show that the morphism of algebraic spaces
$p : Z \times_\mathcal{Y} \mathcal{X} \to Z$ is quasi-compact.
Let $U \subset Z$ be affine open. Then
$p^{-1}(U) = U \times_\mathcal{Y} Z$
and the algebraic space $U \times_\mathcal{Y} Z$
is quasi-compact by assumption (2). Hence $p$ is quasi-compact, see
Morphisms of Spaces, Lemma \ref{spaces-morphisms-lemma-quasi-compact-local}.
\end{proof}

\end{document}
